% =========================================================
% CHAPTER 2: LITERATURE SURVEY
% =========================================================

\chapter{LITERATURE SURVEY}
\label{chap:literature_survey}

The field of query optimization has a long history, and the application of
machine learning to this domain has accelerated significantly in recent years.
This chapter reviews the most relevant prior work, spanning traditional
optimizer evaluation, learned query optimization, database configuration
tuning, concept drift detection, and uncertainty estimation. Each paper is
discussed in terms of its contribution, its achievements, and how it relates
to the design and evaluation of QueryPilot.

Leis et al. \cite{leis2015good} conducted a foundational empirical study
titled ``How Good Are Query Optimizers, Really?'' published in the Proceedings
of the VLDB Endowment in 2015. The authors systematically evaluated
cardinality estimation errors across several commercial and open-source
database systems, demonstrating that even small estimation errors compound
across join trees and lead to dramatically suboptimal plan selections. Their
key finding was that cardinality misestimation, rather than the cost model
itself, is the primary source of optimizer failure in practice. This work
establishes the foundational motivation for external optimization frameworks
like QueryPilot: since internal statistics can become unrepresentative,
learning from actual execution feedback provides a more reliable basis for
optimization decisions than relying solely on pre-computed statistics.

Marcus and Papaemmanouil \cite{marcus2019neo} introduced Neo, a learned query
optimizer published in the Proceedings of the VLDB Endowment in 2019. Neo uses
a value network trained on query execution feedback to evaluate candidate
execution plans produced by the database's own optimizer, selecting the plan
with the highest estimated value. The system achieves significant improvements
over the native optimizer on the Join Order Benchmark by learning to avoid the
compound cardinality estimation errors identified by Leis et al. However, Neo
is tightly integrated with the database's plan enumeration engine and requires
access to the plan search space. Crucially, Neo does not address what should
happen when the model encounters queries that are dissimilar from its training
distribution. QueryPilot differs from Neo by operating as an external
middleware without requiring plan enumeration access, and by explicitly
handling model uncertainty and workload drift.

Marcus et al. \cite{marcus2021bao} presented Bao (Bandit Optimizer) at ACM
SIGMOD 2021, framing optimizer hint selection as a contextual bandit problem.
Bao uses a tree-convolutional neural network over the query plan structure to
recommend one of a small, fixed set of hint configurations for each query.
Bao's formulation is particularly relevant to QueryPilot because both systems
work within a predefined hint space rather than enumerating all possible plans.
Bao demonstrates strong performance on TPC-H and the Join Order Benchmark and
has been extended to handle non-stationary workloads using a sliding window
approach. QueryPilot differs from Bao in that it adds an explicit uncertainty
estimation layer using ensemble disagreement and a drift detection mechanism
using the Page-Hinkley test, and uses a conservative fallback decision when
uncertainty is high rather than always applying the bandit's recommendation.

Van Aken et al. \cite{vanaken2017ottertune} introduced OtterTune, a system for
automatic database management system tuning through large-scale machine
learning, presented at ACM SIGMOD in 2017. OtterTune uses Gaussian process
regression to recommend database configuration knob settings --- such as buffer
pool size, thread count, and I/O scheduler parameters --- based on a repository
of previously observed workload-configuration-performance tuples. Unlike
query-level optimizers, OtterTune operates at the configuration level and
targets long-term system throughput rather than per-query execution time.
QueryPilot operates at the per-query level and does not modify database
configuration. OtterTune's use of uncertainty-aware regression via Gaussian
processes is conceptually related to QueryPilot's ensemble uncertainty
estimation, but the two systems address fundamentally different optimization
targets and operate at different granularities.

Gama et al. \cite{gama2014survey} published a comprehensive survey titled ``A
Survey on Concept Drift Adaptation'' in ACM Computing Surveys in 2014,
reviewing the broad landscape of methods for learning from non-stationary data
streams. The survey distinguishes between abrupt drift, gradual drift,
recurring drift, and incremental drift, and categorizes detection methods as
statistical process control methods, windowing methods, and ensemble methods.
This taxonomy directly informed QueryPilot's experimental design, which
evaluates performance under four scenarios corresponding to stable, sudden
shift, gradual shift, and data skew workloads. The survey identifies
statistical process control methods such as the Page-Hinkley test as well
suited for detecting abrupt changes in data streams, which supports
QueryPilot's choice of drift detector.

Bifet and Gavalda \cite{bifet2007adwin} introduced the ADWIN (Adaptive
Windowing) algorithm at the 2007 SIAM International Conference on Data Mining.
ADWIN maintains a variable-length window over a data stream and automatically
shrinks the window when a statistical change is detected, ensuring that the
active window always contains data from the current distribution. ADWIN
provides formal guarantees on false positive and false negative rates and is
widely used in stream mining applications. QueryPilot uses the Page-Hinkley
test instead of ADWIN, primarily because Page-Hinkley is computationally
lighter for high-frequency query streams and does not require storing the full
window of observations. Both methods are evaluated in the literature as
effective for detecting abrupt shifts in data streams.

Hinkley \cite{hinkley1971inference} originally proposed the sequential
change-point detection method known as the Page-Hinkley test in the journal
Biometrika in 1971. The test maintains a cumulative sum statistic that
increases when observed values exceed the running mean by more than a small
tolerance, and signals a change when this statistic exceeds a threshold. The
test is well-suited to query execution time streams because it is
computationally efficient with constant memory requirements, requires no window
storage, and is sensitive to sustained upward shifts in mean execution time.
QueryPilot implements the Page-Hinkley test with parameters $\delta = 0.005$
and $\lambda = 25$, maintaining one detector per unique query signature in
addition to a global detector.

Lakshminarayanan et al. \cite{lakshminarayanan2017ensembles} demonstrated in
their NeurIPS 2017 paper ``Simple and Scalable Predictive Uncertainty
Estimation Using Deep Ensembles'' that training multiple independently
initialized neural networks and using their output disagreement as an
uncertainty signal is a practical and effective approach to uncertainty
quantification. While the original paper focuses on deep neural networks,
the ensemble disagreement principle generalizes to any ensemble of predictive
models. QueryPilot applies this principle to an ensemble of three scikit-learn
models --- Linear Regression, Random Forest Regressor, and Gradient Boosting
Regressor --- where the standard deviation of the three predictions divided by
their mean serves as the relative uncertainty estimate. This approach requires
no Bayesian inference and introduces negligible computational overhead compared
to training a single model.
